\addtocontents{toc}{\protect\Needspace{3\baselineskip}}
\originalchapter{8.821 课程详讲：世界面作用量、约束与光锥规范}
本章对应第 10 讲和第 11 讲的经典及正则准备. 同一世界面方程可以描述开弦或闭弦, 但边界、零模归一化和独立振子数不同. 本章以 $\hbar=c=1$、负时间号为约定, 目标空间维数写 $D$, 世界面指标 $a,b=0,1$, 横向指标 $I=2,\ldots,D-1$.

\originalsection{从弦面积读取张力和波速}
目标度规为 $G_{\mu\nu}(X)$, 嵌入 $X^\mu(\tau,\sigma)$ 给诱导度规
\[
h_{ab}=G_{\mu\nu}(X)\partial_aX^\mu\partial_bX^\nu.
\]
对时向世界面, $\det h<0$, 面积作用量取
\[
S_{\rm NG}=-T_0\int\dd\tau\dd\sigma\sqrt{-\det h},
\qquad T_0=\frac1{2\pi\alpha'},\quad \ell_s=\sqrt{\alpha'}.
\]
前面的负号使静止弦的能量为正; 与不同文献比较时须同时比较度规和作用量约定.

取近静止直弦, 用静态规范 $X^0=\tau,X^1=\sigma$, 横向位移记作 $Y^I$. 则
\[
h_{00}=-1+\dot Y^2,\quad h_{11}=1+Y'^2,
\quad h_{01}=\dot Y\cdot Y',
\]
以及 $-\det h=(1-\dot Y^2)(1+Y'^2)+(\dot Y\cdot Y')^2$. 对小斜率、小横向速度展开平方根,
\[
\mathcal L=-T_0+\frac{T_0}2(\dot Y^2-Y'^2)+O((\partial Y)^4).
\]
第一项给单位长度能量 $T_0$, 二次项给 $\ddot Y-Y''=0$. 速度为 1, 恢复单位就是光速. 非相对论弦的波速 $\sqrt{T/\mu}$ 只有在张力与质量密度的相对论关系下才变成 $c$, 不能把所有日常绳波都说成光速波.

\originalsection{Polyakov 的辅助度规如何被消去}
引入独立世界面度规 $\gamma_{ab}$,
\begin{equation}\label{eq:Polyakov}
S_P=-\frac{T_0}2\int\dd^2\sigma\sqrt{-\gamma}
\gamma^{ab}G_{\mu\nu}(X)\partial_aX^\mu\partial_bX^\nu.
\end{equation}
不同于 $h_{ab}$, $\gamma_{ab}$ 暂时不是嵌入计算出来的量. 对其逆度规变分, 用 $\delta\sqrt{-\gamma}=-(\sqrt{-\gamma}/2)\gamma_{ab}\delta\gamma^{ab}$, 得约束
\[
h_{ab}-\frac12\gamma_{ab}\gamma^{cd}h_{cd}=0.
\]
对非退化世界面, 这要求 $\gamma_{ab}=\Omega^2h_{ab}$. 在二维 $\sqrt{-\gamma}=\Omega^2\sqrt{-h}$、$\gamma^{ab}=\Omega^{-2}h^{ab}$, 且 $h^{ab}h_{ab}=2$, 代回恰好得到 $S_{\rm NG}$. 因而经典方程在相应条件下等价. 在量子理论里还需规定积分测度、规范固定、规范固定鬼场与 Weyl 一致性; 不能只用这个代入操作就证明所有量子路径积分也已等价.

这个作用量有目标 Poincaré 全局对称（平直 $G=\eta$ 时）, 世界面坐标重参数化和局部 Weyl 变换. 重参数化下 $X$ 是世界面标量, 而度规满足
\[
\gamma'_{ab}(\sigma')=
\frac{\partial\sigma^c}{\partial\sigma'^a}
\frac{\partial\sigma^d}{\partial\sigma'^b}\gamma_{cd}(\sigma).
\]
Weyl 变换 $\gamma_{ab}\to\ee^{2\omega}\gamma_{ab}$ 时, 二维行列式平方根乘 $\ee^{2\omega}$, 逆度规乘 $\ee^{-2\omega}$, 两者抵消. 这是一项局部冗余, 不等于对目标空间坐标同时作缩放.

\originalsection{Euler 项和模参数不能被遗漏}
在欧氏世界面上, 常数膨胀子的拓扑项是
\[
S_\chi=\phi_0\left[
\frac1{4\pi}\int_\Sigma\sqrt\gamma R^{(2)}\dd^2\sigma
+\frac1{2\pi}\int_{\partial\Sigma}K\dd s\right]
=\phi_0\chi(\Sigma).
\]
这里假定光滑边界; 有角点时 Gauss--Bonnet 还需角点项. 对定向面, $\chi=2-2h-b$, $h$ 为亏格、$b$ 为真正边界数. Weyl 变换给
\[
R'^{(2)}=\ee^{-2\omega}(R^{(2)}-2\nabla^2\omega),
\quad K'=\ee^{-\omega}(K+n^a\partial_a\omega),
\quad \dd s'=\ee^\omega\dd s.
\]
体积分多出的 $-2\nabla^2\omega$ 经散度定理, 与边界 $n^a\partial_a\omega$ 正好抵消. 因此是完整体加边界的 Euler 项保持不变, 单独漏掉边界项会破坏这项检查.

二维曲率恒等式 $R_{ab}=\tfrac12\gamma_{ab}R$ 使纯 Einstein 项的局部度规方程平凡, 但不使整个弦量子理论平凡. Weyl 异常和全局模参数仍重要. 以 $g_s=\ee^{\phi_0}$ 定义, 闭弦把手使权重增加 $g_s^2$, 盘式开放过程相应有 $g_o^2\propto g_s$. 耦合的归一化系数还需与具体振幅和场归一化匹配.

局部可以把二维度规取为共形平直, 再用 Weyl 取 $\gamma_{ab}=\eta_{ab}$. 全局并不能消去所有形状参数. 欧氏环面例如用 $z\sim z+1\sim z+\tau_{\rm mod}$, $\operatorname{Im}\tau_{\rm mod}>0$, 描述相对边粘合. 这个复参数改变周期方向的相对长度与斜度, 是不能被局部平直化删除的模. 等价粘合存在模变换识别. 第 10 讲用环带粘合参数 $w$ 表示同一构造, 可联系为 $w=\ee^{2\pi\ii\tau_{\rm mod}}$.

两把手曲面可再去掉并粘合一对圆盘, 其一般光滑形状有三个复模参数; 更高亏格 $h\ge2$ 有 $3h-3$ 个复模. 这些参数在散射路径积分中需要积分. 它们说明“取平直规范”是局部规范处理, 不是把所有曲面都等同于一个毫无形状信息的平面. 完整模空间测度和超模参数的构造超出本讲经典谱准备, 不在这里假装已经完成.

\originalsection{边界项决定开放弦的端点}
平直目标中对 $X^\mu$ 变分, 体方程是
\[
\partial_a(\sqrt{-\gamma}\gamma^{ab}\partial_bX^\mu)=0,
\]
而边界项正比于 $\int_{\partial\Sigma}\dd s\,
\delta X_\mu n^a\partial_aX^\mu$. 在共形规范、空间边界 $\sigma=0,\pi$, 可使它消失的两类基本选择为
\[
\text{Neumann: }X'^\mu=0,\qquad
\text{Dirichlet: }\delta X^\mu=0.
\]
前者端点可在该方向运动, 后者把端点值固定在所给子流形上. 同一个开弦可在不同目标方向分别采用两类, 导出 D 膜结构. 它们来自变分原理的一致性, 不是为了“计算方便”任意丢掉端点通量.

闭弦没有这些端点, 取 $\sigma\sim\sigma+2\pi$, 嵌入及有关度规数据按适当周期条件粘合. 紧致目标方向还可有缠绕, 后面 8.251 的紧致弦课程将单独处理; 这里先用非紧致平直目标和零缠绕来推自由谱.

\originalsection{波动方程之外还存在两个约束}
在 $\gamma=\eta$ 下, 体方程变成 $\ddot X^\mu-X''^\mu=0$, 但度规变分的方程仍必须满足:
\[
\dot X^2+X'^2=0,\qquad \dot X\cdot X'=0.
\]
若只把平直规范作用量当作 $D$ 个互不受约束标量去量子化, 会把负号时间场的非物理态一并留下. 先变分再规范固定, 保留 Virasoro 约束, 是区分弦与一组普通二维自由标量的必要步骤.

波动方程的通解可写 $X^\mu=x^\mu+v^\mu\tau+X_L^\mu(\tau+\sigma)+X_R^\mu(\tau-\sigma)$. 闭弦左右周期函数独立. Neumann 开弦两端的导数条件将左右函数联系成反射后同一周期函数, 最终形成驻波 $\cos n\sigma$. 等式允许吸收进 $x^\mu$ 的常数差, 无需把常数零模重复数一次.

取 $\sigma^\pm=\tau\pm\sigma$, 共形规范仍留下独立变换 $\sigma^+\to f(\sigma^+)$、$\sigma^-\to g(\sigma^-)$ 及补偿 Weyl. 在 $p^+\ne0$、没有妨碍此规范的全局退化下, 可令
\[
X^+=\frac{X^0+X^1}{\sqrt2}=x^++v^+\tau.
\]
目标线元为 $-2\dd X^+\dd X^-+\dd X^I\dd X^I$. 两个约束于是给
\[
2v^+\dot X^-=(\dot X^I)^2+(X'^I)^2,
\qquad v^+X'^- =\dot X^I X'^I.
\]
故 $X^-$ 的导数由横向场确定, 仅余一个位置常数零模. 横向 $D-2$ 个场是独立振荡自由度, 它们具有正号动能. 对闭弦还要确保积分一圈后 $X^-$ 正确周期, 得到下一章的层级匹配.

\begin{example}
若某个 Neumann 横向场含单模 $X^2=A\cos\tau\cos\sigma$, 波动方程成立, $\sigma=0,\pi$ 的 $X'^2$ 也为零. 但不能因此把 $X^-=v^-\tau$ 任意选成常数斜率, 因为 $\dot X^2X'^2=A^2\sin\tau\cos\tau\sin\sigma\cos\sigma$ 一般非零, 约束要求相应的 $X'^-$ 振荡. 这个例子显示纵向场的受迫变化如何把一个看似普通站波嵌入相对论弦.
\end{example}
